% SLT Quasi-Singular Analysis Paper
\documentclass[11pt]{article}

% Essential packages for research papers
\usepackage[utf8]{inputenc}
\usepackage[T1]{fontenc}
\usepackage{amsmath,amsfonts,amssymb}
\usepackage{graphicx}
\usepackage{hyperref}
\usepackage{natbib}
\usepackage{booktabs}
\usepackage{algorithm}
\usepackage{algorithmic}
\usepackage{geometry}

% Page layout
\geometry{margin=1in}

% Title and authors
\title{Posterior Concentration in Normal Mixture Models: An SLT Perspective}
\author{Your Name\footnote{Affiliation and email}}
\date{\today}

\begin{document}
\maketitle

\begin{abstract}
Write your abstract here. This template demonstrates how to integrate LaTeX papers with the research monorepo workflow.
\end{abstract}

\section{Introduction}

This is an example paper that demonstrates:
\begin{itemize}
\item How to reference figures from your notebooks
\item Mathematical notation: $E = mc^2$
\item Code integration with the monorepo packages
\end{itemize}

\section{Methods}

You can reference your monorepo code and results:

\begin{equation}
\mathcal{L}(\theta) = \sum_{i=1}^{n} \log p(x_i | \theta)
\end{equation}

\section{Results}

Include figures generated from your notebooks in \texttt{projects/}.

\begin{figure}[h]
\centering
% \includegraphics[width=0.7\textwidth]{figures/posterior_concentration.pdf}
\caption{Posterior concentration across sample sizes. Generated from \texttt{fit\_dataset.ipynb}.}
\label{fig:posterior}
\end{figure}

\section{Conclusion}

Your conclusions here.

% Bibliography
\bibliographystyle{plainnat}
\bibliography{references}

\end{document}
